\chapter{Diagrammi dei package e di sequenza}

Il diagramma delle classi del Capitolo 7 descrive gli elementi del sistema uno per uno. Per sistemi reali servono altri due punti
di vista: come le classi sono \textbf{raggruppate} in moduli (il diagramma dei \textbf{package}) e come gli oggetti
\textbf{collaborano nel tempo} per realizzare una funzionalità (il diagramma di \textbf{sequenza}).

% ============================================================
\section{Il diagramma dei package}

\dfn{Diagramma dei package}{
  Diagramma diverso da quello delle classi, anche se ne condivide molti aspetti ed elementi. Vi compaiono:
  \begin{itemize}
    \item i \textbf{package}, disegnati come cartelle con una linguetta;
    \item le relazioni di \textbf{dipendenza} tra i package;
    \item le \textbf{interfacce} e le relazioni di \textbf{realizzazione}.
  \end{itemize}
}

Talvolta diagramma dei package e diagramma delle classi si fondono in un unico diagramma: dentro ogni package si disegna il
diagramma delle classi che gli appartengono.

\Needspace{8\baselineskip}
I diagrammi dei package si usano di solito come diagrammi di \textbf{dettaglio}, che mostrano la reale implementazione di un
sistema:
\begin{itemize}
  \item servono quando si sceglie l'\textbf{architettura} del progetto;
  \item la fase successiva è il disegno dei diagrammi delle classi dell'implementazione, uno per ogni package.
\end{itemize}

\begin{esempio}{L'architettura dei progetti del corso}
\begin{center}
\begin{tikzpicture}
  \umlpkg{gui}{(0,3.4)}{5.6}{1.3}{GUI}
  \node[font=\scriptsize, text width=5.2cm, align=center] at ($(gui.center)+(0,-0.15)$) {Qui mettiamo tutte le classi relative ai frame dell'interfaccia grafica};
  \umlpkg{ctr}{(0,1.4)}{5.6}{1.3}{Controller}
  \node[font=\scriptsize, text width=5.2cm, align=center] at ($(ctr.center)+(0,-0.15)$) {Nei nostri progetti, per semplicità, il controller avrà una sola classe};
  \umlpkg{mod}{(0,-0.6)}{5.6}{1.3}{Model}
  \node[font=\scriptsize, text width=5.2cm, align=center] at ($(mod.center)+(0,-0.15)$) {Qui riproduciamo il diagramma delle classi progettato in fase di analisi del dominio};
  \draw[dep] (gui.south) -- (ctr.north);
  \draw[dep] (ctr.south) -- (mod.north);
\end{tikzpicture}
\end{center}
L'interfaccia grafica dipende dal controller, che a sua volta dipende dal modello. Il modello non dipende da nessuno: si può
cambiare l'interfaccia senza toccare le classi del dominio.
\end{esempio}

\begin{esempio}{Package e classi insieme}
\begin{center}
\begin{tikzpicture}
  % GUI
  \umlpkg{gui}{(0,3.6)}{5.4}{1.6}{GUI}
  \node[classe1=1.6cm] (home) at (1.2,4.2) {Home};
  \node[classe1=1.8cm] (ns) at (4.0,4.2) {NuovaSocieta};
  \draw[dep] (home) -- (ns);
  % Controller
  \umlpkg{ctrp}{(0,1.2)}{5.4}{1.8}{Controller}
  \umlclass[4.0cm]{ctr}{(2.7,1.8)}{Controller}{}{+ nuovaSocieta(nome, prezzo)}
  \draw[dep] (ns.south) -- (ns.south |- ctr.north);
  % Model
  \umlpkg{mod}{(0,-1.7)}{7.4}{1.7}{Model}
  \node[classe1=1.4cm] (bo) at (0.9,-1.05) {Borsa};
  \umlclass[1.9cm]{li}{(3.4,-1.0)}{Listino}{}{+ addSocieta()}
  \node[classe1=1.5cm] (so) at (6.3,-1.05) {Societa};
  \draw[line width=0.7pt] (bo) -- node[molt, below, pos=0.15]{1} node[molt, below, pos=0.8]{1} (li.west |- bo);
  \draw[line width=0.7pt] (li.east |- so) -- node[molt, below, pos=0.15]{1} node[molt, below, pos=0.8]{1..*} (so);
  \draw[dep] ($(ctr.south)+(-1.2,0)$) -- (bo.north);
  \draw[dep] ($(ctr.south)+(1.2,0)$) -- node[stereo, right]{«create»} (so.north);
  % DAO
  \umlpkg{daop}{(6.1,1.2)}{3.2}{1.8}{DAO}
  \umlclass[2.5cm]{dao}{(7.7,1.75)}{{\tiny\normalfont «interface»}\\ListinoDAO}{}{+ addSocietaDB()}
  \draw[dep] (ctr.east |- dao) -- (dao.west);
  % Implementazione
  \umlpkg{impp}{(9.9,1.2)}{4.9}{1.8}{ImplementazionePostgresDAO}
  \node[classe1=4.5cm] (imp) at (12.35,1.75) {ListinoImplementazionePostgresDAO};
  \draw[realizza] (imp.west) -- (dao.east |- imp.west);
  % Database
  \umlpkg{dbp}{(9.9,-1.9)}{4.9}{1.9}{Database}
  \umlclass[3.4cm]{db}{(12.35,-1.2)}{ConnessioneDatabase}{+ connection}{}
  \draw[dep] (imp.south) -- (db.north);
\end{tikzpicture}
\end{center}
\begin{itemize}
  \item Nel package \textbf{GUI} la schermata Home apre la schermata per inserire una nuova società.
  \item L'unica classe del package \textbf{Controller} usa le classi del \textbf{Model} (e crea le società, «create»).
  \item Il Controller non parla direttamente con il database: dipende dall'\textbf{interfaccia} \texttt{ListinoDAO} (\textit{Data
        Access Object}), realizzata da una classe specifica per PostgreSQL che usa la connessione al \textbf{Database}. Per passare a
        un altro database basta scrivere un'altra implementazione dell'interfaccia.
\end{itemize}
\end{esempio}

% ============================================================
\section{I diagrammi di interazione}

\dfn{Diagrammi di interazione}{
  Diagrammi che modellano gli aspetti \textbf{dinamici} di un sistema software, mettendo in evidenza le \textbf{interazioni} tra
  gli elementi che lo compongono (descritti nei diagrammi strutturali). Appartengono alla famiglia dei diagrammi
  \textbf{comportamentali}.
}

\Needspace{8\baselineskip}
Esistono quattro tipi di diagrammi di interazione:
\begin{itemize}
  \item i \textbf{sequence diagram} (diagrammi di sequenza), di gran lunga i più usati;
  \item i \textbf{communication diagram} (in UML 1 si chiamavano \textit{collaboration diagram});
  \item gli \textbf{interaction overview diagram};
  \item i \textbf{timing diagram}.
\end{itemize}

\dfn{Diagramma di sequenza}{
  Modella le interazioni tra uno o più \textbf{attori} e il sistema software durante l'esecuzione di uno \textbf{scenario}. Come
  prevede il paradigma a oggetti, le interazioni tra attori e sistema, e tra le parti del sistema, sono rappresentate come
  \textbf{messaggi}. Il \textbf{tempo scorre dall'alto verso il basso}.
}

\begin{esempio}{Il calcolo del prezzo di un ordine}
\begin{center}
\begin{tikzpicture}
  \foreach \x/\n/\t in {0/o/an Order, 3/ol/an Order Line, 6/pr/aProduct, 9/cu/aCustomer}{
    \node[partec, minimum width=2.2cm] (\n) at (\x,0) {\t};
    \draw[vita] (\n.south) -- (\x,-7.0);
  }
  % messaggio trovato
  \fill (-2.4,-1.0) circle (0.07);
  \draw[sinc] (-2.4,-1.0) -- node[msgl]{calculatePrice} (-0.08,-1.0);
  \attiv{0}{-1.0}{-6.6}
  \draw[sinc] (0.08,-1.5) -- node[msgl]{getQuantity} (2.92,-1.5);
  \attiv{3}{-1.5}{-1.9}
  \draw[sinc] (0.08,-2.2) -- node[msgl]{getProduct} (2.92,-2.2);
  \attiv{3}{-2.2}{-2.7}
  \draw[ritorno] (2.92,-2.7) -- node[msgl]{aProduct} (0.08,-2.7);
  \draw[sinc] (0.08,-3.3) -- node[msgl, pos=0.35]{getPricingDetails} (5.92,-3.3);
  \attiv{6}{-3.3}{-3.8}
  % chiamate a se stesso
  \draw[sinc] (0.08,-4.1) -- (1.4,-4.1) -- (1.4,-4.4) -- (0.24,-4.4); \node[font=\tiny, anchor=west] at (1.45,-4.25) {calculateBasePrice};
  \draw[fill=white, line width=0.6pt] (0.08,-4.4) rectangle (0.24,-4.9);
  \draw[sinc] (0.08,-5.1) -- (1.4,-5.1) -- (1.4,-5.4) -- (0.24,-5.4); \node[font=\tiny, anchor=west] at (1.45,-5.25) {calculateDiscounts};
  \draw[fill=white, line width=0.6pt] (0.08,-5.4) rectangle (0.24,-6.3);
  \draw[sinc] (0.24,-5.8) -- node[msgl, pos=0.6]{getDiscountInfo} (8.92,-5.8);
  \attiv{9}{-5.8}{-6.3}
  % annotazioni
  \begin{scope}[every node/.style={font=\tiny\itshape, text=black!60}, ann/.style={-{Stealth[length=1.4mm]}, dotted, black!55}]
    \node (a1) at (-2.2,-1.9) {messaggio trovato};    \draw[ann] (a1.north) -- (-1.9,-1.1);
    \node (a2) at (4.5,-1.0) {partecipante};          \draw[ann] (a2.north west) -- (3.5,-0.35);
    \node (a3) at (7.4,-1.5) {linea di vita};         \draw[ann] (a3.west) -- (6.05,-1.5);
    \node (a4) at (7.4,-2.9) {attivazione};           \draw[ann] (a4.south west) -- (6.1,-3.4);
    \node (a5) at (4.2,-2.4) {ritorno};               \draw[ann] (a5.west) -- (2.5,-2.65);
    \node (a6) at (4.2,-4.7) {chiamata a se stesso};  \draw[ann] (a6.west) -- (1.45,-4.35);
    \node (a7) at (7.4,-5.2) {messaggio};             \draw[ann] (a7.south) -- (7.2,-5.7);
  \end{scope}
\end{tikzpicture}
\end{center}
L'ordine riceve la richiesta \texttt{calculatePrice} da qualcuno che non compare nel diagramma (\textbf{messaggio trovato}, con il
pallino). Chiede quantità e prodotto alla riga d'ordine, chiede i dettagli del prezzo al prodotto, calcola il prezzo base e gli
sconti chiamando \textbf{se stesso}, e per gli sconti interroga il cliente.
\end{esempio}

\subsection{Gli elementi}

\begin{table}[H]
\centering
\renewcommand{\arraystretch}{1.4}
\begin{tabularx}{\textwidth}{>{\bfseries}P{3.2cm} L}
\toprule
Elemento & Descrizione \\
\midrule
Oggetti (istanze di classi) & Rettangoli con il nome della classe e l'identificatore dell'oggetto sottolineati (notazione UML 1,
ad esempio \underline{\texttt{o1:Order}}), oppure semplicemente un nome da cui si capisce che è un'istanza della classe
(\texttt{anOrder}, \texttt{aProduct}). \\
Attori o endpoint & Riportati a sinistra, con frecce verso gli oggetti del sistema. Rappresentano la persona o il metodo che avvia
l'elaborazione. Si possono omettere se lo scenario è avviato da un altro scenario. \\
Messaggi & Frecce da un attore a un oggetto o tra due oggetti. L'ordine dall'alto in basso è l'ordine in cui vengono scambiati. \\
Linea di vita (\textit{lifeline}) & Linea verticale tratteggiata che parte dal rettangolo dell'oggetto: indica il periodo di vita
dell'oggetto, dalla creazione alla distruzione. \\
Barra di attivazione & Rettangolo sulla linea di vita di un oggetto che riceve un messaggio: rappresenta il tempo necessario per
elaborare la richiesta. \\
\bottomrule
\end{tabularx}
\end{table}

\Needspace{12\baselineskip}
\begin{center}
\begin{tikzpicture}
  \draw[sinc] (0,0) -- node[msgl]{messaggio sincrono} (3.4,0);
  \node[note, anchor=west] at (3.7,0) {freccia \textbf{triangolare piena}: chi invia aspetta la risposta};
  \draw[asinc] (0,-0.7) -- node[msgl]{messaggio asincrono} (3.4,-0.7);
  \node[note, anchor=west] at (3.7,-0.7) {freccia \textbf{semplice}: chi invia prosegue senza aspettare};
  \draw[ritorno] (3.4,-1.4) -- node[msgl]{ritorno} (0,-1.4);
  \node[note, anchor=west] at (3.7,-1.4) {linea \textbf{tratteggiata}: la risposta (si può omettere)};
  \draw[sinc] (1.5,-2.0) -- (2.6,-2.0) -- (2.6,-2.3) -- (1.58,-2.3);
  \draw[fill=white, line width=0.6pt] (1.42,-1.9) rectangle (1.58,-2.6);
  \node[note, anchor=west] at (3.7,-2.2) {un messaggio può restare dentro lo stesso oggetto (\textbf{autoanello})};
\end{tikzpicture}
\end{center}

\subsection{Creazione e distruzione degli oggetti}

\Needspace{10\baselineskip}
\begin{itemize}
  \item Le linee di vita degli oggetti che esistono già (ad esempio quelli \texttt{static}) partono dalla \textbf{cima} del diagramma.
  \item Una chiamata al \textbf{costruttore} crea un oggetto: il rettangolo dell'oggetto si disegna alla \textbf{punta} della freccia
        del messaggio di creazione, più in basso degli altri.
  \item La \textbf{distruzione} di un oggetto si indica con una \textbf{croce} che interrompe la sua linea di vita.
\end{itemize}

\begin{esempio}{Una query al database}
\begin{center}
\begin{tikzpicture}
  \node[partec, minimum width=2.0cm] (h) at (0,0) {a Handler};
  \draw[vita] (h.south) -- (0,-7.2);
  \draw[sinc] (-2.6,-0.9) -- node[msgl]{query database} (-0.08,-0.9);
  \attiv{0}{-0.9}{-6.8}
  % creazione di a Query Command
  \node[partec, minimum width=2.0cm] (q) at (3.4,-1.5) {a Query\\Command};
  \draw[sinc] (0.08,-1.5) -- node[msgl]{new} (q.west);
  \draw[vita] (q.south) -- (3.4,-6.5);
  \attiv{3.4}{-2.0}{-6.4}
  % creazione di a Database Statement
  \node[partec, minimum width=2.0cm] (d) at (6.8,-2.4) {a Database\\Statement};
  \draw[sinc] (3.48,-2.4) -- node[msgl]{new} (d.west);
  \draw[vita] (d.south) -- (6.8,-5.6);
  \draw[sinc] (3.48,-3.3) -- node[msgl]{execute} (6.72,-3.3);
  \attiv{6.8}{-3.3}{-3.8}
  \draw[ritorno] (6.72,-3.8) -- node[msgl]{results} (3.48,-3.8);
  \draw[sinc] (3.48,-4.2) -- (4.6,-4.2) -- (4.6,-4.5) -- (3.64,-4.5); \node[font=\tiny, anchor=west] at (4.65,-4.35) {extract results};
  \draw[fill=white, line width=0.6pt] (3.48,-4.5) rectangle (3.64,-5.0);
  \draw[sinc] (3.48,-5.6) -- node[msgl]{close} (6.72,-5.6);
  \draw[line width=1pt] (6.62,-5.42) -- (6.98,-5.78) (6.62,-5.78) -- (6.98,-5.42);
  \draw[ritorno] (3.4,-6.5) -- node[msgl]{results} (0.08,-6.5);
  \draw[line width=1pt] (3.22,-6.32) -- (3.58,-6.68) (3.22,-6.68) -- (3.58,-6.32);
  \begin{scope}[every node/.style={font=\tiny\itshape, text=black!60}, ann/.style={-{Stealth[length=1.4mm]}, dotted, black!55}]
    \node (c1) at (1.6,-2.6) {creazione};                          \draw[ann] (c1.north) -- (1.6,-1.6);
    \node[align=center] (c2) at (8.6,-4.8) {distruzione\\da parte di\\un altro oggetto}; \draw[ann] (c2.south west) -- (7.0,-5.55);
    \node (c3) at (5.3,-7.0) {autodistruzione};                    \draw[ann] (c3.west) -- (3.65,-6.6);
  \end{scope}
\end{tikzpicture}
\end{center}
Il gestore crea un comando di query, che a sua volta crea uno statement del database, lo esegue, ne estrae i risultati e lo
chiude (lo statement viene distrutto dal comando). Alla fine il comando restituisce i risultati al gestore e si distrugge.
\end{esempio}

\begin{esempio}{L'iscrizione a un corso, con il suo diagramma delle classi}
\begin{center}
\begin{tikzpicture}
  \stickman{(-1.4,-0.3)}{}
  \draw[vita] (-1.4,-0.85) -- (-1.4,-3.6);
  \attiv{-1.4}{-1.0}{-3.3}
  \node[partec, minimum width=2.2cm] (cs) at (2.0,0) {\underline{:CourseSection}};
  \draw[vita] (cs.south) -- (2.0,-3.6);
  \node[partec, minimum width=2.2cm] (st) at (8.4,0) {\underline{:Student}};
  \draw[vita] (st.south) -- (8.4,-3.6);
  \draw[sinc] (-1.32,-1.1) -- node[msgl]{requestToRegister} (1.92,-1.1);
  \attiv{2.0}{-1.1}{-3.2}
  \node[partec, minimum width=2.0cm] (re) at (5.2,-1.6) {\underline{:Registration}};
  \draw[dipend] (2.08,-1.6) -- node[msgl]{«create»} (re.west);
  \draw[vita] (re.south) -- (5.2,-3.6);
  \attiv{5.2}{-1.9}{-3.0}
  \draw[sinc] (5.28,-2.2) -- node[msgl]{addToSchedule} (8.32,-2.2);
  \attiv{8.4}{-2.2}{-2.6}
  \draw[sinc] (5.12,-2.8) -- node[msgl]{addToRegistrationList} (2.08,-2.8);
  % diagramma delle classi
  \umlclass[2.4cm]{c1}{(-0.8,-5.0)}{Course}{}{getPrerequisite}
  \umlclass[3.0cm]{c2}{(2.6,-5.0)}{CourseSection}{}{requestToRegister\\addToRegistrationList}
  \node[classe1=2.2cm] (c3) at (5.9,-4.75) {Registration};
  \umlclass[2.4cm]{c4}{(9.0,-5.0)}{Student}{}{addToSchedule\\hasPassedCourse}
  \draw[line width=0.7pt] (c1.east |- c3) -- node[molt, above, pos=0.8]{*} (c2.west |- c3);
  \draw[line width=0.7pt] (c2.east |- c3) -- node[molt, above, pos=0.8]{*} (c3.west);
  \draw[line width=0.7pt] (c3.east) -- node[molt, above, pos=0.15]{*} (c4.west |- c3);
\end{tikzpicture}
\end{center}
Lo studente chiede alla sezione del corso di iscriversi; la sezione \textbf{crea} una Registration, che aggiunge il corso al piano
di studi dello studente e poi si aggiunge alla lista delle iscrizioni della sezione. Ogni messaggio che un oggetto riceve è
un'\textbf{operazione} della sua classe nel diagramma delle classi: i due diagrammi devono essere \textbf{coerenti}.
\end{esempio}

\subsection{Quando si usa}

\Needspace{10\baselineskip}
In fase di \textbf{progettazione di dettaglio}, i diagrammi di sequenza descrivono la \textbf{realizzazione dei metodi}:
\begin{itemize}
  \item gli oggetti sono istanze delle classi di progetto di dettaglio;
  \item i messaggi sono le chiamate di metodo sugli oggetti;
  \item al posto degli elementi architetturali compaiono gli oggetti delle classi che vi danno accesso (ad esempio, invece del
        «database», l'oggetto JDBC su cui si eseguono le query).
\end{itemize}

\warning{Non per gli algoritmi}{
  Per descrivere \textbf{algoritmi}, i diagrammi di sequenza non sono lo strumento più adatto: di solito si preferiscono i diagrammi
  di attività o di stato (Capitolo 9). Un diagramma di sequenza mostra bene \textbf{chi parla con chi}, non la logica di calcolo.
}

\subsection{Cicli e condizioni}

\Needspace{12\baselineskip}
Cicli e condizioni si indicano con un \textbf{riquadro} (\textit{frame}) che racchiude una sottosequenza di messaggi. Nell'angolo
in alto a sinistra si indica il costrutto; questa notazione è stata introdotta con UML 2.

\begin{table}[H]
\centering
\renewcommand{\arraystretch}{1.4}
\begin{tabularx}{\textwidth}{>{\bfseries\ttfamily}P{1.4cm} P{2.6cm} L}
\toprule
\normalfont\textbf{Costrutto} & \textbf{Corrisponde a} & \textbf{Uso} \\
\midrule
loop & \texttt{while-do}, \texttt{do-while} & Ciclo: la condizione si indica tra parentesi quadre all'inizio o alla fine. \\
alt & \texttt{if-then-else} & La condizione si indica in cima; se ci sono rami \textit{else}, una linea tratteggiata separa la zona
\textit{then} dalla zona \textit{else}, eventualmente con un'altra condizione accanto a \texttt{[else]}. \\
opt & \texttt{if-then} & Racchiude una sottosequenza eseguita solo se la condizione in cima è vera. \\
\bottomrule
\end{tabularx}
\end{table}

Esistono anche altri costrutti, ad esempio per il parallelismo (\texttt{par}) o le regioni critiche (\texttt{critical}).

\begin{esempio}{La spedizione di un ordine}
\noindent\begin{minipage}[c]{0.33\textwidth}
\begin{lstlisting}[language=Java, basicstyle=\scriptsize\ttfamily]
procedure dispatch
  foreach (lineitem)
    if (product.value > 10K)
      careful.dispatch
    else
      regular.dispatch
    end if
  end for
  if (needsConfirmation)
    messenger.confirm
end procedure
\end{lstlisting}
\end{minipage}\hfill
\begin{minipage}[c]{0.64\textwidth}
\centering
\begin{tikzpicture}[scale=0.92, every node/.append style={transform shape}]
  \foreach \x/\n/\t in {0/o/:Order, 2.6/ca/careful :\\Distributor, 4.9/re/regular :\\Distributor, 7.2/me/:Messenger}{
    \node[partec, minimum width=1.9cm] (\n) at (\x,0) {\t};
    \draw[vita] (\n.south) -- (\x,-7.3);
  }
  \fill (-1.4,-0.9) circle (0.07);
  \draw[sinc] (-1.4,-0.9) -- node[msgl]{dispatch} (-0.08,-0.9);
  \attiv{0}{-0.9}{-7.0}
  \sdframe{-1.0}{-1.3}{6.0}{-5.2}{loop}{[for each line item]}
  \sdframe{-0.7}{-1.9}{5.8}{-4.9}{alt}{[value > \$10000]}
  \draw[sinc] (0.08,-2.7) -- node[msgl]{dispatch} (2.52,-2.7);
  \attiv{2.6}{-2.7}{-3.1}
  \draw[dashed, line width=0.6pt] (-0.7,-3.5) -- (5.8,-3.5);
  \node[font=\tiny, anchor=north west] at (-0.6,-3.55) {[else]};
  \draw[sinc] (0.08,-4.2) -- node[msgl]{dispatch} (4.82,-4.2);
  \attiv{4.9}{-4.2}{-4.6}
  \sdframe{-1.0}{-5.5}{7.9}{-6.9}{opt}{[needsConfirmation]}
  \draw[sinc] (0.08,-6.3) -- node[msgl]{confirm} (7.12,-6.3);
  \attiv{7.2}{-6.3}{-6.7}
\end{tikzpicture}
\end{minipage}

\medskip
Per ogni riga dell'ordine (\textbf{loop}), se il valore del prodotto supera i 10.000 dollari la spedizione è affidata al
distributore «attento», altrimenti a quello normale (\textbf{alt}, con il ramo \texttt{[else]} sotto la linea tratteggiata).
Alla fine, solo se serve una conferma (\textbf{opt}), si avvisa il messaggero.
\end{esempio}

\info{{Riquadri sì, ma con moderazione}}{
  Quando l'algoritmo da modellare si fa complesso, è buona norma usare altri tipi di diagramma. Un diagramma di sequenza pieno di
  riquadri annidati è difficile da leggere quanto il codice che descrive.
}

% ============================================================
\section{Esempio: dal codice al diagramma di sequenza}

Immaginiamo di voler disegnare il diagramma di sequenza di una funzionalità \textbf{già implementata}. Sarà inevitabilmente un
diagramma di \textbf{dettaglio}, in cui compaiono tutte le classi realmente presenti nel codice, comprese quelle dell'interfaccia
utente. La funzionalità è \textbf{Aggiungi nuova società (nome, valore azione)} del videogioco di borsa del Capitolo 7, con
l'architettura GUI--Controller--Model vista sopra.

\Needspace{16\baselineskip}
\textbf{1. La schermata Home.} Servono due input, quindi si crea una nuova finestra e la si apre dalla Home:
\begin{lstlisting}[language=Java]
aggiungiNuovaSocietaButton.addActionListener(new ActionListener() {
    @Override
    public void actionPerformed(ActionEvent e) {
        AggiungiNuovaSocieta aggiungiNuovaSocieta = new AggiungiNuovaSocieta(controller, frame);
        aggiungiNuovaSocieta.frame.setVisible(true);
        frame.setVisible(false);
    }
});
\end{lstlisting}

\Needspace{18\baselineskip}
\textbf{2. Il pulsante OK.} Alla pressione bisogna verificare che i dati siano validi (lo si farà in futuro), inserire i dati nel
modello tramite il controller e tornare alla schermata precedente (\textit{Annulla} farà solo quest'ultima operazione):
\begin{lstlisting}[language=Java]
OKButton.addActionListener(new ActionListener() {
    @Override
    public void actionPerformed(ActionEvent e) {
        // aggiunge la societa' al listino
        controller.aggiungiSocietaBorsa(nomeSocietaText.getText(),
                Double.parseDouble(valoreAzioneText.getText()));
        frame.setVisible(false);
        frameChiamante.setVisible(true);
        frame.dispose();
    }
});
\end{lstlisting}

\Needspace{24\baselineskip}
\textbf{3. Controller e modello.} Il controller si limita a inoltrare la richiesta all'oggetto \texttt{borsa}, che ha una lista di
società inizialmente vuota; \texttt{aggiungiSocieta} crea una nuova società e la aggiunge alla lista.
\begin{lstlisting}[language=Java]
// nel Controller
public void aggiungiSocietaBorsa(String text, double parseDouble) {
    borsa.aggiungiSocieta(text, parseDouble);
}

public class Borsa {
    private ArrayList<Societa> societa;
    public Borsa() { societa = new ArrayList<Societa>(); }

    public void aggiungiSocieta(String text, double parseDouble) {
        Societa s = new Societa(text, parseDouble);
        societa.add(s);
    }
}

public class Societa {
    private Double valoreAzione;
    private String nome;
    public Societa(String n, Double valore) { nome = n; valoreAzione = valore; }
}
\end{lstlisting}

\begin{figure}[H]
\centering
\begin{tikzpicture}[scale=0.93, every node/.append style={transform shape}]
  \draw[line width=0.6pt] (-0.9,1.3) rectangle (15.4,-8.4);
  \draw[line width=0.6pt, fill=white] (-0.9,1.3) -- ++(2.6,0) -- ++(0,-0.2) -- ++(-0.12,-0.14) -- (-0.9,0.96) -- cycle;
  \node[font=\tiny, anchor=north west, inner sep=1.5pt] at (-0.85,1.28) {\textbf{sd} Aggiungi nuova società};
  \stickman{(0,-0.1)}{a User}
  \draw[vita] (0,-0.9) -- (0,-8.1);
  \foreach \x/\n/\t in {1.9/ho/a Home, 7.3/ct/a Controller, 10.2/bo/a Borsa, 14.4/al/a ArrayList\\<Societa>}{
    \node[partec, minimum width=1.6cm] (\n) at (\x,-0.2) {\t};
    \draw[vita] (\n.south) -- (\x,-8.1);
  }
  \draw[sinc] (0,-1.3) -- node[msgl]{1: click(Aggiungi)} (1.82,-1.3);
  \attiv{1.9}{-1.3}{-2.1}
  \node[partec, minimum width=2.3cm] (fo) at (4.9,-1.7) {an AggiungiNuovaSocieta};
  \draw[dipend] (1.98,-1.7) -- (fo.west); \node[font=\tiny, anchor=south] at (2.85,-1.68) {2: «create»};
  \draw[vita] (fo.south) -- (4.9,-7.6);
  \draw[sinc] (0,-2.6) -- node[msgl, pos=0.3]{3: setNomeSocieta} (4.82,-2.6);
  \attiv{4.9}{-2.6}{-2.85}
  \draw[sinc] (0,-3.2) -- node[msgl, pos=0.3]{4: setPrezzoAzione} (4.82,-3.2);
  \attiv{4.9}{-3.2}{-3.45}
  \draw[sinc] (0,-3.8) -- node[msgl, pos=0.3]{5: click(OK)} (4.82,-3.8);
  \attiv{4.9}{-3.8}{-7.3}
  \draw[sinc] (4.98,-4.3) -- node[msgl, align=center]{6: aggiungiSocietaBorsa\\(nome, prezzo)} (7.22,-4.3);
  \attiv{7.3}{-4.3}{-6.5}
  \draw[sinc] (7.38,-4.8) -- node[msgl, align=center]{7: aggiungiSocieta\\(nome, prezzo)} (10.12,-4.8);
  \attiv{10.2}{-4.8}{-6.2}
  \node[partec, minimum width=1.5cm] (so) at (12.3,-5.2) {a Societa};
  \draw[dipend] (10.28,-5.2) -- (so.west); \node[font=\tiny, anchor=south] at (10.95,-5.18) {8: «create»};
  \draw[vita] (so.south) -- (12.3,-8.1);
  \draw[sinc] (10.28,-5.9) -- node[msgl, pos=0.75]{9: add(s)} (14.32,-5.9);
  \attiv{14.4}{-5.9}{-6.15}
  \draw[ritorno] (10.12,-6.2) -- node[msgl]{10} (7.38,-6.2);
  \draw[ritorno] (7.22,-6.5) -- node[msgl]{11} (4.98,-6.5);
  \draw[sinc] (4.82,-6.9) -- node[msgl]{12: setVisible(true)} (1.98,-6.9);
  \attiv{1.9}{-6.9}{-7.2}
  \draw[line width=1pt] (4.72,-7.42) -- (5.08,-7.78) (4.72,-7.78) -- (5.08,-7.42);
  \node[font=\tiny, anchor=west] at (5.15,-7.6) {13: dispose()};
\end{tikzpicture}
\caption{Il diagramma di sequenza di «Aggiungi nuova società», ricavato dal codice.}
\end{figure}

\Needspace{6\baselineskip}
Si notino la \textbf{creazione} della finestra (2) e della società (8), disegnate alla punta della freccia; i \textbf{ritorni}
tratteggiati (10 e 11) quando i metodi terminano; e la \textbf{distruzione} della finestra, che si chiude da sola con
\texttt{dispose()}. L'oggetto \texttt{ArrayList} esiste già, perché è stato creato dal costruttore di \texttt{Borsa}.

\info{{Un nome, tre versioni}}{
  Nelle slide lo stesso metodo del controller si chiama \texttt{aggiungiSocietaBorsa} nel codice, \texttt{aggiungiSocietaListino}
  nel diagramma di sequenza e \texttt{nuovaSocieta} nel diagramma dei package. Qui si usa il nome del codice. Mantenere
  \textbf{coerenti} codice e diagrammi è esattamente il lavoro che, nei progetti reali, gli strumenti di co-generazione (come il
  plug-in \textit{Diagrams} di IntelliJ) aiutano a fare.
}

% ============================================================
\section{Esercizi}

\begin{esercizio}{Dal codice al diagramma}
Disegna il diagramma di sequenza della chiamata \texttt{carrello.totale()}.

\Needspace{14\baselineskip}
\begin{lstlisting}[language=Java]
public class Carrello {
    private List<Riga> righe;
    private Cliente cliente;

    public double totale() {
        double t = 0;
        for (Riga r : righe) {
            t += r.getSubtotale();
        }
        if (cliente.isPremium()) {
            t = applicaSconto(t);
        }
        return t;
    }
}
\end{lstlisting}

\textbf{Soluzione.} Il \texttt{for} diventa un riquadro \textbf{loop}, l'\texttt{if} senza \texttt{else} un riquadro \textbf{opt},
e \texttt{applicaSconto} è una chiamata a se stesso.
\begin{center}
\begin{tikzpicture}
  \foreach \x/\n/\t in {0/ca/:Carrello, 3.6/ri/:Riga, 7.2/cl/:Cliente}{
    \node[partec, minimum width=1.8cm] (\n) at (\x,0) {\t};
    \draw[vita] (\n.south) -- (\x,-5.8);
  }
  \fill (-2.0,-0.8) circle (0.07);
  \draw[sinc] (-2.0,-0.8) -- node[msgl]{totale()} (-0.08,-0.8);
  \attiv{0}{-0.8}{-5.4}
  \sdframe{-0.9}{-1.2}{4.6}{-2.6}{loop}{[per ogni riga]}
  \draw[sinc] (0.08,-1.9) -- node[msgl]{getSubtotale()} (3.52,-1.9);
  \attiv{3.6}{-1.9}{-2.3}
  \draw[ritorno] (3.52,-2.3) -- (0.08,-2.3);
  \draw[sinc] (0.08,-3.0) -- node[msgl]{isPremium()} (7.12,-3.0);
  \attiv{7.2}{-3.0}{-3.3}
  \draw[ritorno] (7.12,-3.3) -- (0.08,-3.3);
  \sdframe{-0.9}{-3.6}{4.6}{-5.0}{opt}{[cliente premium]}
  \draw[sinc] (0.08,-4.2) -- (1.4,-4.2) -- (1.4,-4.5) -- (0.24,-4.5); \node[font=\tiny, anchor=west] at (1.45,-4.35) {applicaSconto(t)};
  \draw[fill=white, line width=0.6pt] (0.08,-4.5) rectangle (0.24,-4.8);
  \draw[ritorno] (-0.08,-5.4) -- node[msgl]{t} (-2.0,-5.4);
\end{tikzpicture}
\end{center}
\end{esercizio}

\begin{esercizio}{Vero o falso?}
\begin{enumerate}
  \item In un diagramma di sequenza il tempo scorre da sinistra verso destra.
  \item Un messaggio asincrono si disegna con una freccia a punta triangolare piena.
  \item Il rettangolo di un oggetto creato durante lo scenario si disegna in cima al diagramma, come gli altri.
  \item Un \texttt{if} senza \texttt{else} si rappresenta con un riquadro \texttt{opt}.
  \item Nel diagramma dei package, una freccia tratteggiata tra due package indica una dipendenza.
\end{enumerate}

\textbf{Soluzione.} (1) Falso: dall'alto verso il basso. (2) Falso: la punta piena è per i messaggi sincroni; gli asincroni hanno
una freccia semplice. (3) Falso: si disegna alla punta della freccia del messaggio di creazione. (4) Vero. (5) Vero.
\end{esercizio}

\gbox{In sintesi}{azure-gradient-6!80!black}{
\begin{itemize}
  \item Il \textbf{diagramma dei package} mostra i package (cartelle), le dipendenze tra loro e le interfacce; si usa per scegliere
        l'architettura, ad esempio GUI--Controller--Model con i DAO per l'accesso ai dati. Può contenere i diagrammi delle classi.
  \item I \textbf{diagrammi di interazione} sono quattro: di sequenza, di comunicazione, interaction overview e timing.
  \item Il \textbf{diagramma di sequenza} mostra i \textbf{messaggi} scambiati da attori e oggetti in uno scenario, dall'alto verso
        il basso: oggetti, linee di vita, barre di attivazione, messaggi sincroni (punta piena), asincroni (punta semplice) e di
        ritorno (tratteggiati), chiamate a se stessi, creazione e distruzione (croce).
  \item I riquadri \textbf{loop}, \textbf{alt} e \textbf{opt} (UML 2) rappresentano cicli e condizioni, ma per gli algoritmi
        complessi sono meglio i diagrammi di attività o di stato.
  \item Nel progetto di dettaglio i messaggi sono chiamate di metodo, e il diagramma di sequenza deve essere \textbf{coerente} con
        il diagramma delle classi e con il codice.
\end{itemize}
}
